
\documentclass{article}
\usepackage{colortbl}
\usepackage{makecell}
\usepackage{multirow}
\usepackage{supertabular}

\begin{document}

\newcounter{utterance}

\centering \large Interaction Transcript for game `matchit\_ascii', experiment `same\_grid', episode 2 with claude{-}opus{-}4{-}8{-}azure.
\vspace{24pt}

{ \footnotesize  \setcounter{utterance}{1}
\setlength{\tabcolsep}{0pt}
\begin{supertabular}{c@{$\;$}|p{.15\linewidth}@{}p{.15\linewidth}p{.15\linewidth}p{.15\linewidth}p{.15\linewidth}p{.15\linewidth}}
    \# & $\;$A & \multicolumn{4}{c}{Game Master} & $\;\:$B\\
    \hline

    \theutterance \stepcounter{utterance}  
    & & \multicolumn{4}{p{0.6\linewidth}}{
        \cellcolor[rgb]{0.9,0.9,0.9}{
            \makecell[{{p{\linewidth}}}]{
                \texttt{\tiny{[P1$\langle$GM]}}
                \texttt{Dalyvaujate bendrame spėjimo žaidime.} \\
\texttt{Tikslas {-} išsiaiškinti, ar šis tinklelis ir kitas paveikslėlis, kurį matau tik aš, yra vienodi.} \\
\texttt{Pirmiausia apibūdinkite savo tinklelį. Tada aš pateiksiu savo aprašymą ir galėsime vienas kitam užduoti klausimų apie tinklelius, kad išsiaiškintume, ar jie vienodi. Dabar pradėkite savo trumpą tinklelio aprašymą žodžiais "APRAŠYMAS:", po kurių rašykite aprašymą. Nieko daugiau nepridėkite. Tai yra jūsų tinklelis:} \\
\texttt{X ▢ ▢ ▢ X} \\
\texttt{X X ▢ X X} \\
\texttt{X ▢ X ▢ X} \\
\texttt{X ▢ ▢ ▢ X} \\
\texttt{X ▢ ▢ ▢ X} \\
            }
        }
    }
    & \\ \\

    \theutterance \stepcounter{utterance}  
    & \multicolumn{4}{p{0.6\linewidth}}{
        \cellcolor[rgb]{0.9,0.9,0.9}{
            \makecell[{{p{\linewidth}}}]{
                \texttt{\tiny{[P1$\rangle$GM]}}
                \texttt{APRAŠYMAS: Tai 5x5 tinklelis. Pirmoje eilutėje X yra kairiajame ir dešiniajame kampuose (1 ir 5 stulpeliai), viduryje tuščia. Antroje eilutėje X yra 1, 2, 4 ir 5 stulpeliuose, viduryje (3 stulpelis) tuščia. Trečioje eilutėje X yra 1, 3 ir 5 stulpeliuose. Ketvirta ir penkta eilutės vienodos – X tik 1 ir 5 stulpeliuose (kraštuose), vidurys tuščias. Visas kairysis ir dešinysis stulpeliai pilnai užpildyti X.} \\
            }
        }
    }
    & & \\ \\

    \theutterance \stepcounter{utterance}  
    & & & \multicolumn{2}{p{0.3\linewidth}}{
        \cellcolor[rgb]{0.9,0.9,0.9}{
            \makecell[{{p{\linewidth}}}]{
                \texttt{\tiny{[GM$|$GM]}}
                \texttt{continue} \\
            }
        }
    }
    & & \\ \\

    \theutterance \stepcounter{utterance}  
    & & \multicolumn{4}{p{0.6\linewidth}}{
        \cellcolor[rgb]{0.9,0.9,0.9}{
            \makecell[{{p{\linewidth}}}]{
                \texttt{\tiny{[GM$\rangle$P2]}}
                \texttt{Dalyvaujate bendrame spėjimo žaidime.} \\
\texttt{Tikslas {-} išsiaiškinti, ar šis tinklelis ir kitas paveikslėlis, kurį matau tik aš, yra vienodi.} \\
\texttt{Pirmiausia apibūdinkite savo tinklelį. Tada aš pateiksiu savo aprašymą ir galėsime vienas kitam užduoti klausimų apie tinklelius, kad išsiaiškintume, ar jie vienodi. Dabar pradėkite savo trumpą tinklelio aprašymą žodžiais "APRAŠYMAS:", po kurių rašykite aprašymą. Nieko daugiau nepridėkite. Tai yra jūsų tinklelis:} \\
\texttt{X ▢ ▢ ▢ X} \\
\texttt{X X ▢ X X} \\
\texttt{X ▢ X ▢ X} \\
\texttt{X ▢ ▢ ▢ X} \\
\texttt{X ▢ ▢ ▢ X} \\
            }
        }
    }
    & \\ \\

    \theutterance \stepcounter{utterance}  
    & & & \multicolumn{4}{p{0.6\linewidth}}{
        \cellcolor[rgb]{0.9,0.9,0.9}{
            \makecell[{{p{\linewidth}}}]{
                \texttt{\tiny{[GM$\langle$P2]}}
                \texttt{APRAŠYMAS: Tai 5x5 tinklelis. 1 eilutė: X, tuščia, tuščia, tuščia, X. 2 eilutė: X, X, tuščia, X, X. 3 eilutė: X, tuščia, X, tuščia, X. 4 eilutė: X, tuščia, tuščia, tuščia, X. 5 eilutė: X, tuščia, tuščia, tuščia, X. Kraštinės kairysis ir dešinysis stulpeliai visiškai užpildyti X. Viduryje yra "X" formos raštas.} \\
            }
        }
    }
     \\ \\

    \theutterance \stepcounter{utterance}  
    & & & \multicolumn{2}{p{0.3\linewidth}}{
        \cellcolor[rgb]{0.9,0.9,0.9}{
            \makecell[{{p{\linewidth}}}]{
                \texttt{\tiny{[GM$|$GM]}}
                \texttt{continue} \\
            }
        }
    }
    & & \\ \\

    \theutterance \stepcounter{utterance}  
    & & \multicolumn{4}{p{0.6\linewidth}}{
        \cellcolor[rgb]{0.9,0.9,0.9}{
            \makecell[{{p{\linewidth}}}]{
                \texttt{\tiny{[GM$\rangle$P2]}}
                \texttt{Tai mano APRAŠYMASAPRAŠYMAS: Tai 5x5 tinklelis. Pirmoje eilutėje X yra kairiajame ir dešiniajame kampuose (1 ir 5 stulpeliai), viduryje tuščia. Antroje eilutėje X yra 1, 2, 4 ir 5 stulpeliuose, viduryje (3 stulpelis) tuščia. Trečioje eilutėje X yra 1, 3 ir 5 stulpeliuose. Ketvirta ir penkta eilutės vienodos – X tik 1 ir 5 stulpeliuose (kraštuose), vidurys tuščias. Visas kairysis ir dešinysis stulpeliai pilnai užpildyti X.} \\
\texttt{Dabar užduokite klausimą, kad sužinotumėte naujus mano tinklelio aspektus, kurie gali skirtis nuo jūsų tinklelio. Pradėkite "KLAUSIMAS:" ir nieko daugiau nepridėkite.} \\
            }
        }
    }
    & \\ \\

    \theutterance \stepcounter{utterance}  
    & & & \multicolumn{4}{p{0.6\linewidth}}{
        \cellcolor[rgb]{0.9,0.9,0.9}{
            \makecell[{{p{\linewidth}}}]{
                \texttt{\tiny{[GM$\langle$P2]}}
                \texttt{KLAUSIMAS: Ar jūsų tinklelyje trečioje eilutėje 2 ir 4 stulpeliai yra tušti (t. y. X tik 1, 3 ir 5 stulpeliuose)?} \\
            }
        }
    }
     \\ \\

    \theutterance \stepcounter{utterance}  
    & & & \multicolumn{2}{p{0.3\linewidth}}{
        \cellcolor[rgb]{0.9,0.9,0.9}{
            \makecell[{{p{\linewidth}}}]{
                \texttt{\tiny{[GM$|$GM]}}
                \texttt{continue} \\
            }
        }
    }
    & & \\ \\

    \theutterance \stepcounter{utterance}  
    & & \multicolumn{4}{p{0.6\linewidth}}{
        \cellcolor[rgb]{0.9,0.9,0.9}{
            \makecell[{{p{\linewidth}}}]{
                \texttt{\tiny{[P1$\langle$GM]}}
                \texttt{Tai mano APRAŠYMASAPRAŠYMAS: Tai 5x5 tinklelis. 1 eilutė: X, tuščia, tuščia, tuščia, X. 2 eilutė: X, X, tuščia, X, X. 3 eilutė: X, tuščia, X, tuščia, X. 4 eilutė: X, tuščia, tuščia, tuščia, X. 5 eilutė: X, tuščia, tuščia, tuščia, X. Kraštinės kairysis ir dešinysis stulpeliai visiškai užpildyti X. Viduryje yra "X" formos raštas.} \\
\texttt{KLAUSIMAS: Ar jūsų tinklelyje trečioje eilutėje 2 ir 4 stulpeliai yra tušti (t. y. X tik 1, 3 ir 5 stulpeliuose)?Pradėkite atsakymą „ATSAKYMAS:" ir nieko daugiau nepridėkite.} \\
            }
        }
    }
    & \\ \\

    \theutterance \stepcounter{utterance}  
    & \multicolumn{4}{p{0.6\linewidth}}{
        \cellcolor[rgb]{0.9,0.9,0.9}{
            \makecell[{{p{\linewidth}}}]{
                \texttt{\tiny{[P1$\rangle$GM]}}
                \texttt{ATSAKYMAS: Taip, mano tinklelyje trečioje eilutėje 2 ir 4 stulpeliai yra tušti – X yra tik 1, 3 ir 5 stulpeliuose. Mūsų aprašymai visiškai sutampa, todėl tinkleliai yra vienodi.} \\
            }
        }
    }
    & & \\ \\

    \theutterance \stepcounter{utterance}  
    & & & \multicolumn{2}{p{0.3\linewidth}}{
        \cellcolor[rgb]{0.9,0.9,0.9}{
            \makecell[{{p{\linewidth}}}]{
                \texttt{\tiny{[GM$|$GM]}}
                \texttt{continue} \\
            }
        }
    }
    & & \\ \\

    \theutterance \stepcounter{utterance}  
    & & \multicolumn{4}{p{0.6\linewidth}}{
        \cellcolor[rgb]{0.9,0.9,0.9}{
            \makecell[{{p{\linewidth}}}]{
                \texttt{\tiny{[P1$\langle$GM]}}
                \texttt{Dabar užduokite klausimą, kad sužinotumėte naujus mano tinklelio aspektus, kurie gali skirtis nuo jūsų tinklelio. Pradėkite "KLAUSIMAS:" ir nieko daugiau nepridėkite.} \\
            }
        }
    }
    & \\ \\

    \theutterance \stepcounter{utterance}  
    & \multicolumn{4}{p{0.6\linewidth}}{
        \cellcolor[rgb]{0.9,0.9,0.9}{
            \makecell[{{p{\linewidth}}}]{
                \texttt{\tiny{[P1$\rangle$GM]}}
                \texttt{KLAUSIMAS: Ar jūsų tinklelyje antroje eilutėje X yra 1, 2, 4 ir 5 stulpeliuose, o vidurinis (3 stulpelis) tuščias?} \\
            }
        }
    }
    & & \\ \\

    \theutterance \stepcounter{utterance}  
    & & & \multicolumn{2}{p{0.3\linewidth}}{
        \cellcolor[rgb]{0.9,0.9,0.9}{
            \makecell[{{p{\linewidth}}}]{
                \texttt{\tiny{[GM$|$GM]}}
                \texttt{continue} \\
            }
        }
    }
    & & \\ \\

    \theutterance \stepcounter{utterance}  
    & & \multicolumn{4}{p{0.6\linewidth}}{
        \cellcolor[rgb]{0.9,0.9,0.9}{
            \makecell[{{p{\linewidth}}}]{
                \texttt{\tiny{[GM$\rangle$P2]}}
                \texttt{ATSAKYMAS: Taip, mano tinklelyje trečioje eilutėje 2 ir 4 stulpeliai yra tušti – X yra tik 1, 3 ir 5 stulpeliuose. Mūsų aprašymai visiškai sutampa, todėl tinkleliai yra vienodi.} \\
\texttt{KLAUSIMAS: Ar jūsų tinklelyje antroje eilutėje X yra 1, 2, 4 ir 5 stulpeliuose, o vidurinis (3 stulpelis) tuščias?Pradėkite atsakymą „ATSAKYMAS:" ir nieko daugiau nepridėkite.} \\
            }
        }
    }
    & \\ \\

    \theutterance \stepcounter{utterance}  
    & & & \multicolumn{4}{p{0.6\linewidth}}{
        \cellcolor[rgb]{0.9,0.9,0.9}{
            \makecell[{{p{\linewidth}}}]{
                \texttt{\tiny{[GM$\langle$P2]}}
                \texttt{ATSAKYMAS: Taip, mano tinklelyje antroje eilutėje X yra 1, 2, 4 ir 5 stulpeliuose, o 3 stulpelis tuščias. Visi aprašymai sutampa, todėl tinkleliai yra vienodi.} \\
            }
        }
    }
     \\ \\

    \theutterance \stepcounter{utterance}  
    & & & \multicolumn{2}{p{0.3\linewidth}}{
        \cellcolor[rgb]{0.9,0.9,0.9}{
            \makecell[{{p{\linewidth}}}]{
                \texttt{\tiny{[GM$|$GM]}}
                \texttt{continue} \\
            }
        }
    }
    & & \\ \\

    \theutterance \stepcounter{utterance}  
    & & \multicolumn{4}{p{0.6\linewidth}}{
        \cellcolor[rgb]{0.9,0.9,0.9}{
            \makecell[{{p{\linewidth}}}]{
                \texttt{\tiny{[GM$\rangle$P2]}}
                \texttt{Dabar užduokite klausimą, kad sužinotumėte naujus mano tinklelio aspektus, kurie gali skirtis nuo jūsų tinklelio. Pradėkite "KLAUSIMAS:" ir nieko daugiau nepridėkite.} \\
            }
        }
    }
    & \\ \\

    \theutterance \stepcounter{utterance}  
    & & & \multicolumn{4}{p{0.6\linewidth}}{
        \cellcolor[rgb]{0.9,0.9,0.9}{
            \makecell[{{p{\linewidth}}}]{
                \texttt{\tiny{[GM$\langle$P2]}}
                \texttt{KLAUSIMAS: Ar jūsų tinklelyje pirmoje eilutėje X yra tik 1 ir 5 stulpeliuose, o 2, 3 ir 4 stulpeliai tušti?} \\
            }
        }
    }
     \\ \\

    \theutterance \stepcounter{utterance}  
    & & & \multicolumn{2}{p{0.3\linewidth}}{
        \cellcolor[rgb]{0.9,0.9,0.9}{
            \makecell[{{p{\linewidth}}}]{
                \texttt{\tiny{[GM$|$GM]}}
                \texttt{continue} \\
            }
        }
    }
    & & \\ \\

    \theutterance \stepcounter{utterance}  
    & & \multicolumn{4}{p{0.6\linewidth}}{
        \cellcolor[rgb]{0.9,0.9,0.9}{
            \makecell[{{p{\linewidth}}}]{
                \texttt{\tiny{[P1$\langle$GM]}}
                \texttt{ATSAKYMAS: Taip, mano tinklelyje antroje eilutėje X yra 1, 2, 4 ir 5 stulpeliuose, o 3 stulpelis tuščias. Visi aprašymai sutampa, todėl tinkleliai yra vienodi.} \\
\texttt{KLAUSIMAS: Ar jūsų tinklelyje pirmoje eilutėje X yra tik 1 ir 5 stulpeliuose, o 2, 3 ir 4 stulpeliai tušti?Pradėkite atsakymą „ATSAKYMAS:" ir nieko daugiau nepridėkite.} \\
            }
        }
    }
    & \\ \\

    \theutterance \stepcounter{utterance}  
    & \multicolumn{4}{p{0.6\linewidth}}{
        \cellcolor[rgb]{0.9,0.9,0.9}{
            \makecell[{{p{\linewidth}}}]{
                \texttt{\tiny{[P1$\rangle$GM]}}
                \texttt{ATSAKYMAS: Taip, mano tinklelyje pirmoje eilutėje X yra tik 1 ir 5 stulpeliuose, o 2, 3 ir 4 stulpeliai tušti. Visi aprašymai sutampa, todėl tinkleliai yra vienodi.} \\
            }
        }
    }
    & & \\ \\

    \theutterance \stepcounter{utterance}  
    & & & \multicolumn{2}{p{0.3\linewidth}}{
        \cellcolor[rgb]{0.9,0.9,0.9}{
            \makecell[{{p{\linewidth}}}]{
                \texttt{\tiny{[GM$|$GM]}}
                \texttt{continue} \\
            }
        }
    }
    & & \\ \\

    \theutterance \stepcounter{utterance}  
    & & \multicolumn{4}{p{0.6\linewidth}}{
        \cellcolor[rgb]{0.9,0.9,0.9}{
            \makecell[{{p{\linewidth}}}]{
                \texttt{\tiny{[P1$\langle$GM]}}
                \texttt{Dabar užduokite klausimą, kad sužinotumėte naujus mano tinklelio aspektus, kurie gali skirtis nuo jūsų tinklelio. Pradėkite "KLAUSIMAS:" ir nieko daugiau nepridėkite.} \\
            }
        }
    }
    & \\ \\

    \theutterance \stepcounter{utterance}  
    & \multicolumn{4}{p{0.6\linewidth}}{
        \cellcolor[rgb]{0.9,0.9,0.9}{
            \makecell[{{p{\linewidth}}}]{
                \texttt{\tiny{[P1$\rangle$GM]}}
                \texttt{KLAUSIMAS: Ar jūsų tinklelyje ketvirta ir penkta eilutės yra vienodos – X tik 1 ir 5 stulpeliuose, o 2, 3 ir 4 stulpeliai tušti?} \\
            }
        }
    }
    & & \\ \\

    \theutterance \stepcounter{utterance}  
    & & & \multicolumn{2}{p{0.3\linewidth}}{
        \cellcolor[rgb]{0.9,0.9,0.9}{
            \makecell[{{p{\linewidth}}}]{
                \texttt{\tiny{[GM$|$GM]}}
                \texttt{continue} \\
            }
        }
    }
    & & \\ \\

    \theutterance \stepcounter{utterance}  
    & & \multicolumn{4}{p{0.6\linewidth}}{
        \cellcolor[rgb]{0.9,0.9,0.9}{
            \makecell[{{p{\linewidth}}}]{
                \texttt{\tiny{[GM$\rangle$P2]}}
                \texttt{ATSAKYMAS: Taip, mano tinklelyje pirmoje eilutėje X yra tik 1 ir 5 stulpeliuose, o 2, 3 ir 4 stulpeliai tušti. Visi aprašymai sutampa, todėl tinkleliai yra vienodi.} \\
\texttt{KLAUSIMAS: Ar jūsų tinklelyje ketvirta ir penkta eilutės yra vienodos – X tik 1 ir 5 stulpeliuose, o 2, 3 ir 4 stulpeliai tušti?Pradėkite atsakymą „ATSAKYMAS:" ir nieko daugiau nepridėkite.} \\
            }
        }
    }
    & \\ \\

    \theutterance \stepcounter{utterance}  
    & & & \multicolumn{4}{p{0.6\linewidth}}{
        \cellcolor[rgb]{0.9,0.9,0.9}{
            \makecell[{{p{\linewidth}}}]{
                \texttt{\tiny{[GM$\langle$P2]}}
                \texttt{ATSAKYMAS: Taip, mano tinklelyje ketvirta ir penkta eilutės yra vienodos – X tik 1 ir 5 stulpeliuose, o 2, 3 ir 4 stulpeliai tušti. Patikrinę visas eilutes, matome, kad mūsų tinkleliai yra visiškai vienodi.} \\
            }
        }
    }
     \\ \\

    \theutterance \stepcounter{utterance}  
    & & & \multicolumn{2}{p{0.3\linewidth}}{
        \cellcolor[rgb]{0.9,0.9,0.9}{
            \makecell[{{p{\linewidth}}}]{
                \texttt{\tiny{[GM$|$GM]}}
                \texttt{continue} \\
            }
        }
    }
    & & \\ \\

    \theutterance \stepcounter{utterance}  
    & & \multicolumn{4}{p{0.6\linewidth}}{
        \cellcolor[rgb]{0.9,0.9,0.9}{
            \makecell[{{p{\linewidth}}}]{
                \texttt{\tiny{[GM$\rangle$P2]}}
                \texttt{Dabar užduokite klausimą, kad sužinotumėte naujus mano tinklelio aspektus, kurie gali skirtis nuo jūsų tinklelio. Pradėkite "KLAUSIMAS:" ir nieko daugiau nepridėkite.} \\
            }
        }
    }
    & \\ \\

    \theutterance \stepcounter{utterance}  
    & & & \multicolumn{4}{p{0.6\linewidth}}{
        \cellcolor[rgb]{0.9,0.9,0.9}{
            \makecell[{{p{\linewidth}}}]{
                \texttt{\tiny{[GM$\langle$P2]}}
                \texttt{KLAUSIMAS: Ar jūsų tinklelyje iš viso yra 15 X ženklų (patikrinimui: 2+4+3+2+2)?} \\
            }
        }
    }
     \\ \\

    \theutterance \stepcounter{utterance}  
    & & & \multicolumn{2}{p{0.3\linewidth}}{
        \cellcolor[rgb]{0.9,0.9,0.9}{
            \makecell[{{p{\linewidth}}}]{
                \texttt{\tiny{[GM$|$GM]}}
                \texttt{continue} \\
            }
        }
    }
    & & \\ \\

    \theutterance \stepcounter{utterance}  
    & & \multicolumn{4}{p{0.6\linewidth}}{
        \cellcolor[rgb]{0.9,0.9,0.9}{
            \makecell[{{p{\linewidth}}}]{
                \texttt{\tiny{[P1$\langle$GM]}}
                \texttt{ATSAKYMAS: Taip, mano tinklelyje ketvirta ir penkta eilutės yra vienodos – X tik 1 ir 5 stulpeliuose, o 2, 3 ir 4 stulpeliai tušti. Patikrinę visas eilutes, matome, kad mūsų tinkleliai yra visiškai vienodi.} \\
\texttt{KLAUSIMAS: Ar jūsų tinklelyje iš viso yra 15 X ženklų (patikrinimui: 2+4+3+2+2)?Pradėkite atsakymą „ATSAKYMAS:" ir nieko daugiau nepridėkite.} \\
            }
        }
    }
    & \\ \\

    \theutterance \stepcounter{utterance}  
    & \multicolumn{4}{p{0.6\linewidth}}{
        \cellcolor[rgb]{0.9,0.9,0.9}{
            \makecell[{{p{\linewidth}}}]{
                \texttt{\tiny{[P1$\rangle$GM]}}
                \texttt{ATSAKYMAS: Ne, mano tinklelyje iš viso yra 13 X ženklų (2+4+3+2+2 = 13, o ne 15). Tačiau eilutėse nurodyti skaičiai sutampa su jūsų aprašymu, todėl tinkleliai vis tiek yra vienodi – tik jūsų pateikta suma buvo klaidinga.} \\
            }
        }
    }
    & & \\ \\

    \theutterance \stepcounter{utterance}  
    & & & \multicolumn{2}{p{0.3\linewidth}}{
        \cellcolor[rgb]{0.9,0.9,0.9}{
            \makecell[{{p{\linewidth}}}]{
                \texttt{\tiny{[GM$|$GM]}}
                \texttt{continue} \\
            }
        }
    }
    & & \\ \\

    \theutterance \stepcounter{utterance}  
    & & \multicolumn{4}{p{0.6\linewidth}}{
        \cellcolor[rgb]{0.9,0.9,0.9}{
            \makecell[{{p{\linewidth}}}]{
                \texttt{\tiny{[P1$\langle$GM]}}
                \texttt{Dabar užduokite klausimą, kad sužinotumėte naujus mano tinklelio aspektus, kurie gali skirtis nuo jūsų tinklelio. Pradėkite "KLAUSIMAS:" ir nieko daugiau nepridėkite.} \\
            }
        }
    }
    & \\ \\

    \theutterance \stepcounter{utterance}  
    & \multicolumn{4}{p{0.6\linewidth}}{
        \cellcolor[rgb]{0.9,0.9,0.9}{
            \makecell[{{p{\linewidth}}}]{
                \texttt{\tiny{[P1$\rangle$GM]}}
                \texttt{KLAUSIMAS: Ar jūsų tinklelyje visas vidurinis stulpelis (3 stulpelis) yra tuščias, išskyrus trečią eilutę, kurioje yra X?} \\
            }
        }
    }
    & & \\ \\

    \theutterance \stepcounter{utterance}  
    & & & \multicolumn{2}{p{0.3\linewidth}}{
        \cellcolor[rgb]{0.9,0.9,0.9}{
            \makecell[{{p{\linewidth}}}]{
                \texttt{\tiny{[GM$|$GM]}}
                \texttt{continue} \\
            }
        }
    }
    & & \\ \\

    \theutterance \stepcounter{utterance}  
    & & \multicolumn{4}{p{0.6\linewidth}}{
        \cellcolor[rgb]{0.9,0.9,0.9}{
            \makecell[{{p{\linewidth}}}]{
                \texttt{\tiny{[GM$\rangle$P2]}}
                \texttt{ATSAKYMAS: Ne, mano tinklelyje iš viso yra 13 X ženklų (2+4+3+2+2 = 13, o ne 15). Tačiau eilutėse nurodyti skaičiai sutampa su jūsų aprašymu, todėl tinkleliai vis tiek yra vienodi – tik jūsų pateikta suma buvo klaidinga.} \\
\texttt{KLAUSIMAS: Ar jūsų tinklelyje visas vidurinis stulpelis (3 stulpelis) yra tuščias, išskyrus trečią eilutę, kurioje yra X?Pradėkite atsakymą „ATSAKYMAS:" ir nieko daugiau nepridėkite.} \\
            }
        }
    }
    & \\ \\

    \theutterance \stepcounter{utterance}  
    & & & \multicolumn{4}{p{0.6\linewidth}}{
        \cellcolor[rgb]{0.9,0.9,0.9}{
            \makecell[{{p{\linewidth}}}]{
                \texttt{\tiny{[GM$\langle$P2]}}
                \texttt{ATSAKYMAS: Taip, mano tinklelyje visas vidurinis stulpelis (3 stulpelis) yra tuščias, išskyrus trečią eilutę, kurioje yra X. Be to, teisinga, kad iš viso yra 13 X ženklų – mano ankstesnis skaičiavimas buvo klaidingas. Visi eilučių aprašymai sutampa, todėl tinkleliai yra vienodi.} \\
            }
        }
    }
     \\ \\

    \theutterance \stepcounter{utterance}  
    & & & \multicolumn{2}{p{0.3\linewidth}}{
        \cellcolor[rgb]{0.9,0.9,0.9}{
            \makecell[{{p{\linewidth}}}]{
                \texttt{\tiny{[GM$|$GM]}}
                \texttt{continue} \\
            }
        }
    }
    & & \\ \\

    \theutterance \stepcounter{utterance}  
    & & \multicolumn{4}{p{0.6\linewidth}}{
        \cellcolor[rgb]{0.9,0.9,0.9}{
            \makecell[{{p{\linewidth}}}]{
                \texttt{\tiny{[GM$\rangle$P2]}}
                \texttt{Dabar priimkite sprendimą. Kaip manote: ar jūsų tinklelis ir kitas aprašytas tinklelis yra tokie patys? Parašykite "SPRENDIMAS: tas pats tinklelis", jei manote, kad tai tas pats tinklelis, arba "SPRENDIMAS: skirtingi tinkleliai", jei manote, kad tai skirtingi tinkleliai. Nieko daugiau nepridėkite.} \\
            }
        }
    }
    & \\ \\

    \theutterance \stepcounter{utterance}  
    & & & \multicolumn{4}{p{0.6\linewidth}}{
        \cellcolor[rgb]{0.9,0.9,0.9}{
            \makecell[{{p{\linewidth}}}]{
                \texttt{\tiny{[GM$\langle$P2]}}
                \texttt{SPRENDIMAS: tas pats tinklelis} \\
            }
        }
    }
     \\ \\

    \theutterance \stepcounter{utterance}  
    & & & \multicolumn{2}{p{0.3\linewidth}}{
        \cellcolor[rgb]{0.9,0.9,0.9}{
            \makecell[{{p{\linewidth}}}]{
                \texttt{\tiny{[GM$|$GM]}}
                \texttt{continue} \\
            }
        }
    }
    & & \\ \\

    \theutterance \stepcounter{utterance}  
    & & & \multicolumn{2}{p{0.3\linewidth}}{
        \cellcolor[rgb]{0.9,0.9,0.9}{
            \makecell[{{p{\linewidth}}}]{
                \texttt{\tiny{[GM$|$GM]}}
                \texttt{success} \\
            }
        }
    }
    & & \\ \\

    \theutterance \stepcounter{utterance}  
    & & \multicolumn{4}{p{0.6\linewidth}}{
        \cellcolor[rgb]{0.9,0.9,0.9}{
            \makecell[{{p{\linewidth}}}]{
                \texttt{\tiny{[P1$\langle$GM]}}
                \texttt{ATSAKYMAS: Taip, mano tinklelyje visas vidurinis stulpelis (3 stulpelis) yra tuščias, išskyrus trečią eilutę, kurioje yra X. Be to, teisinga, kad iš viso yra 13 X ženklų – mano ankstesnis skaičiavimas buvo klaidingas. Visi eilučių aprašymai sutampa, todėl tinkleliai yra vienodi.} \\
\texttt{Dabar priimkite sprendimą. Kaip manote: ar jūsų tinklelis ir kitas aprašytas tinklelis yra tokie patys? Parašykite "SPRENDIMAS: tas pats tinklelis", jei manote, kad tai tas pats tinklelis, arba "SPRENDIMAS: skirtingi tinkleliai", jei manote, kad tai skirtingi tinkleliai. Nieko daugiau nepridėkite.} \\
            }
        }
    }
    & \\ \\

    \theutterance \stepcounter{utterance}  
    & \multicolumn{4}{p{0.6\linewidth}}{
        \cellcolor[rgb]{0.9,0.9,0.9}{
            \makecell[{{p{\linewidth}}}]{
                \texttt{\tiny{[P1$\rangle$GM]}}
                \texttt{SPRENDIMAS: tas pats tinklelis} \\
            }
        }
    }
    & & \\ \\

    \theutterance \stepcounter{utterance}  
    & & & \multicolumn{2}{p{0.3\linewidth}}{
        \cellcolor[rgb]{0.9,0.9,0.9}{
            \makecell[{{p{\linewidth}}}]{
                \texttt{\tiny{[GM$|$GM]}}
                \texttt{continue} \\
            }
        }
    }
    & & \\ \\

    \theutterance \stepcounter{utterance}  
    & & & \multicolumn{2}{p{0.3\linewidth}}{
        \cellcolor[rgb]{0.9,0.9,0.9}{
            \makecell[{{p{\linewidth}}}]{
                \texttt{\tiny{[GM$|$GM]}}
                \texttt{success} \\
            }
        }
    }
    & & \\ \\

\end{supertabular}
}

\end{document}
